\input _template

\setlanguage{EN}

\Title{Algebra}

\Subtitle{Systems of Equations}

\MathcompsLink{systems-of-equations}

\Author{Patrik Bak}

\sec Introduction

Systems of equations are a common school topic, where we learn many common methods. However, there are many interesting problems where these methods lead to a very complicated procedure or fail completely; yet these problems have very elegant unexpected solutions. Let's take for example
$$
\align
x^2+1 &= 2y, \\
y^2+1 &= 2x.
\endalign
$$
The school approach for this system is to use the \uv{substitution method}. However, when we express $y$ from the first equation and substitute it into the second, we get a fourth-degree equation after a series of adjustments. This can be solved by guessing roots and subsequent polynomial division, but the procedure is tedious\fnote{It is interesting that while general formulas exist for finding roots of equations up to the fourth degree (even though they are very complex), for equations of the fifth and higher degree it was proven in the 19th century (Abel, Galois) that such general formulas do not exist.}. In this material, however, we will show three other more elegant and generally useful methods for solving this example.

\sec Theory 

There are many possible procedures that can be chosen when solving systems of equations. In this material, we will focus on three very frequent techniques:

\begitems \style n
\i Subtracting equations
\i Adding equations
\i Ordering variables
\enditems

Very often we will build upon the techniques of \textit{factorization} or \textit{completing the square}, which we practiced in the previous material.

Before we start solving, one important practical piece of advice: The check -- it is always good to do it and write in the solution that you have done it. Often it is much faster than \textit{precisely} justifying that it is not needed.

\secc Subtracting equations

The most common method for more interesting systems consists of subtracting equations and subsequently factoring. Let's return to the example from the introduction:

\Example{}{
    Solve the system in the set of real numbers using the method of subtracting equations:
    $$
    \align
    x^2+1 &= 2y, \\
    y^2+1 &= 2x. \\
    \endalign
    $$
}{
    Subtracting the second equation from the first, we get
    $$
    \align
    x^2+1-(y^2+1)&=2y-2x,\\
    (x-y)(x+y)&=-2(x-y),\\
    (x-y)(x+y+2)&=0.
    \endalign
    $$
    Thus, either $x=y$ or $x+y=-2$.
    \begitems
    \i If $x=y$, then from $x^2+1=2x$ it follows that $(x-1)^2=0$, that is, $x=y=1$.
    \i If $x+y=-2$, then by substituting into $x^2+1=2y$ we get $x^2+1=2(-2-x)$, which simplifies to $x^2+2x+5=0$. We easily find that this equation has no solution in real numbers.
    \enditems
    The system therefore has a single solution $(x,y)=(1,1)$, of which we can easily convince ourselves by checking.
}

The situation is more complex when we have more equations, then we often need to subtract more pairs of equations and we often cannot avoid case analysis:

\Example{}{
    Solve the system in the set of real numbers using the method of subtracting equations:
    $$
    \align
    x^2-1 &= y+z, \\
    y^2-1 &= z+x, \\
    z^2-1 &= x+y. \\
    \endalign
    $$
}{
    Subtracting the first and second equation, we get
    $$
    x^2-y^2=(y+z)-(z+x)=y-x.
    $$
    After moving to one side, we have
    $$
    \align 
    x^2-y^2 + (x-y) &= 0, \\
    (x-y)(x+y) + (x-y) &= 0, \\
    (x-y)(x+y+1) &= 0.
    \endalign
    $$
    Cyclically, we also obtain
    $$
    \align
    (y-z)(y+z+1)&=0,\\
    (z-x)(z+x+1)&=0.
    \endalign
    $$
    Thus, for every pair of variables it holds: either they are equal, or their sum is $-1$.

    \begitems
    \i If $x=y=z=t$, then from any equation $t^2-1=2t$, i.e., $t^2-2t-1=0$, which gives $t=1\pm\sqrt2$.
    \i If $x=y\neq z$, then from the equations it follows $x+z=-1$. By substituting $y=x$ and $z=-1-x$ into the third equation:
    $$
    (-1-x)^2-1= 2x
    $$
    which simplifies to $x^2=0$, so $x=0$ and $z=-1$. We thus have the solution $(0,0,-1)$. Cyclically in cases $y=z\neq x$ and $z=x \neq y$ we get solutions $(-1,0,0)$ and $(0,-1,0)$.
    \i If no pair were equal, then from $x\neq y$ and $x \neq z$ we have $x+y=-1$ and $x+z=-1$. By comparison we have $x+y=x+z$, so $y=z$, which is a contradiction with the assumption.
    \enditems

    The solutions are the triples $(1+\sqrt2,1+\sqrt2,1+\sqrt2)$, $(1-\sqrt2,1-\sqrt2,1-\sqrt2)$, $(-1,0,0)$, $(0,-1,0)$ and $(0,0,-1)$, of which we can easily convince ourselves by checking.
}

You can practice subtracting equations on the following examples:

\Exercise{}{
    Solve the system of equations in the domain of real numbers:
    $$
    \align
    x^2 + xy &= x + 1,\\
    y^2 + xy &= y + 1.
    \endalign
    $$
}{
    Subtracting the equations we get $x^2-y^2=x-y$, which we adjust to
    $$
    (x-y)(x+y-1)=0.
    $$
    This gives us two possibilities:

    \begitems
    \i If $x=y$, by substituting into the first equation we have $2x^2-x-1=0$. This equation has roots $x=1$ and $x=-1/2$. We thus obtain two solutions: $(1,1)$ and $(-1/2, -1/2)$. The check works.

    \i If $y=1-x$, the first equation transforms to $x^2+x(1-x)=x+1$, from which after adjustment we get $x=x+1$, thus $0=1$. This case leads to no solution.
    \enditems

    The system has exactly two solutions: $(1,1)$ and $(-1/2, -1/2)$.
}

\Exercise{}{
    Solve the system of equations in the domain of real numbers:
    $$
    \align
    xy+x+y &= 3, \\
    yz+y+z &= 3, \\
    zx+z+x &= 3. \\
    \endalign
    $$
}{
    Subtracting the second equation from the first, we get 
    $$
    (xy+x+y) - (yz+y+z) = 0,
    $$
    which we adjust to
    $$
    \align
    y(x-z) + (x-z) &= 0, \\
    (x-z)(y+1) &=0. \\
    \endalign
    $$
    Similarly, we also get 
    $$
    \align 
    (y-x)(y+1) &= 0, \\
    (z-y)(z+1) &= 0. \\
    \endalign
    $$
    If any variable is equal to $-1$, for example $y$, then in the first equation we have a contradiction $-1=3$. Similarly, the remaining variables cannot be equal to $-1$ either. Necessarily, all three are therefore equal.

    Substituting $x=y=z=t$ into the first equation we obtain $t^2+2t-3=0$, which is a quadratic equation with solutions $1$ and $-3$. 

    We can convince ourselves by checking that both found solutions $(1,1,1)$ and $(-3,-3,-3)$ satisfy the system.
}

\Exercise{58th year, home round B, problem 2}{
    Determine all triples $(x,y,z)$ of real numbers for which it holds
    $$
    \align
    x^2+xy&=y^2+z^2,\\
    z^2+zy&=y^2+x^2. \\
    \endalign
    $$
}{
    Subtracting the second equation from the first we get
    $$
    \align
    (x^2+xy)-(z^2+zy)&=(y^2+z^2)-(y^2+x^2),\\
    2(x^2-z^2)+y(x-z)&=0,\\
    2(x-z)(x+z)+y(x-z)&=0,\\
    (x-z)\(2(x+z)+y\)&=0.
    \endalign
    $$
    From there, either $x=z$, or $y=-2(x+z)$.

    \begitems 
    \i If $x=z$, then from the first equation $xy=y^2$, i.e., $y(y-x)=0$. If $y=0$, any $(x,0,x)$ satisfies. If $y=x$, we get $(x,x,x)$ for any real $x$.
    \i If $y=-2(x+z)$, substituting into the first equation we get 
    $$
    \align
    x^2-2x(x+z)&=(-2(x+z))^2+z^2, \\
    x^2-2x^2-2xz &= 4x^2+8xz+4z^2+z^2, \\
    -5x^2 - 10xz - 5z^2 &= 0, \\
    -5(x+z)^2 &= 0.
    \endalign
    $$
    From there $z=-x$ and $y=0$, so the solutions are $(x,0,-x)$ for any real $x$.
    \enditems

    All solutions are exactly the triples $(t,t,t)$, $(t,0,t)$ and $(t,0,-t)$ for any real $t$, of which we can easily convince ourselves by checking.
}

\secc Adding equations

Besides subtracting, it is sometimes a good idea to add the equations as well. There are examples where the resulting equation either suggests something or gives up immediately. Let's return to the example from the introduction:

\Example{}{
    Solve the system in the set of real numbers using the method of adding equations:
    $$
    \align
    x^2+1 &= 2y, \\
    y^2+1 &= 2x. \\
    \endalign
    $$
}{
    Adding the equations from the problem statement, we get
    $$
    x^2+1+y^2+1 = 2y+2x.
    $$
    At first glance, we haven't helped ourselves. However, it is necessary to notice that after moving all terms to one side, we have nothing more than a sum of squares:
    $$
    \align
    x^2+1+y^2+1 &= 2y+2x, \\
    (x^2-2x+1) + (y^2-2y+1) &=0, \\
    (x-1)^2 + (y-1)^2 &= 0.
    \endalign
    $$
    From that we immediately have that it must hold $x=1$ and $y=1$. By checking, we convince ourselves that the pair $(1,1)$ is indeed a solution.
}


We will show an even more complex example with actually three variables:

\Example{}{
    Solve the system in the set of real numbers:
    $$
    \align
    x^2 &= y(2z-x), \\
    y^2 &= z(2x-y), \\
    z^2 &= x(2y-z). \\
    \endalign
    $$
}{
    After adjusting equations to the form
    $$
    x^2+xy=2yz,\qquad y^2+yz=2zx,\qquad z^2+zx=2xy
    $$
    and their subsequent addition we get $x^2+y^2+z^2+xy+yz+zx = 2(xy+yz+zx)$, which is equivalent to
    $$
    x^2+y^2+z^2-xy-yz-zx=0.
    $$
    We can multiply this equality by two and rewrite it as a sum of squares:
    $$
    (x-y)^2+(y-z)^2+(z-x)^2=0.
    $$
    Since we are summing three non-negative numbers, the equality holds only if all three are zero, thus $x-y=0$, $y-z=0$ and $z-x=0$. From this, it follows that $x=y=z$.
    
    By substituting into the original equations, we easily verify that every triple $(t,t,t)$ for $t\in\R$ is a solution.
}

\Exercise{}{
    In the domain of real numbers, solve the system of equations:
    $$
    \align
    a^2+b&=c, \\
    b^2+c&=a, \\
    c^2+a&=b.
    \endalign
    $$
}{
    Adding the given three equations we have
    $$
    a^2+b^2+c^2+a+b+c = a+b+c,
    $$
    from which after adjustment we get
    $$
    a^2+b^2+c^2 = 0.
    $$
    From that we necessarily have $a^2=0$, $b^2=0$ and $c^2=0$, which means $a=b=c=0$.
    
    By substituting into the system, we verify that $(0,0,0)$ is indeed a solution.
}

\Exercise{}{
    In the domain of real numbers, solve the system
    $$
    \align
    z^2 &= 4(x-1),\\
    x^2 &= 4(y-1),\\
    y^2 &= 4(z-1).\\
    \endalign
    $$
}{
    Adding the equations we get
    $$
    x^2+y^2+z^2=4(x+y+z-3).
    $$
    Now we move everything to the left side, write $12$ as $4+4+4$ and complete the square:
    $$
    \align
    &\hphantom{=}x^2-4x+y^2-4y+z^2-4z+12\\
     &=(x^2-4x+4)+(y^2-4y+4)+(z^2-4z+4)\\
     &=(x-2)^2+(y-2)^2+(z-2)^2.
    \endalign
    $$
    This sum should be equal to zero, necessarily therefore each of the three non-negative addends is equal to $0$. By checking we verify that the triple $(2,2,2)$ also satisfies the original system.
}

\Exercise{}{
    Solve the system in real numbers:
    $$
    \align
    x(x+y+z)&=20,\\
    y(x+y+z)&=30,\\
    z(x+y+z)&=50.\\
    \endalign
    $$
}{
    Let's notice that in all three equations a common factor $x+y+z$ occurs, thus by adding the equations we will be able to factor it out:
    $$
    \align
    x(x+y+z)+y(x+y+z)+z(x+y+z) &= 20+30+50, \\
    (x+y+z)(x+y+z) &= 100, \\
    (x+y+z)^2&=100. \\
    \endalign
    $$
    From there we get two possibilities for the value of the sum $x+y+z$:
    $$
    x+y+z = 10 \quad \text{or} \quad x+y+z = -10.
    $$
    We will analyze both cases.
    \begitems
    \i If $x+y+z=10$, substituting into the original equations we get:
    $$
    \align
    x \cdot 10 &= 20 \quad\text{from which}\quad x=2, \\
    y \cdot 10 &= 30 \quad\text{from which}\quad y=3, \\
    z \cdot 10 &= 50 \quad\text{from which}\quad z=5.
    \endalign
    $$
    We obtained the solution $(2,3,5)$.

    \i If $x+y+z=-10$, similarly $(-2,-3,-5)$.
    \enditems
    By checking we convince ourselves that both triples $(2,3,5)$ and $(-2,-3,-5)$ satisfy the system.
}

\secc Ordering variables

Seemingly the trickiest method is to use symmetry and \uv{without loss of generality} order the variables and use estimates.

Without much talk, we will show this on our example:

\Example{}{
    Solve the system in the set of real numbers using the method of ordering variables:
    $$
    \align
    x^2+1 &= 2y, \\
    y^2+1 &= 2x. \\
    \endalign
    $$
}{
    Let's realize that the system is \textit{symmetric} in variables $x$ and $y$, since by swapping variables $x$ and $y$ we get the same problem. In other words, the pair $(x,y)$ is a solution if and only if the pair $(y,x)$ is also. Thus, it suffices to find those pairs where $x \ge y$, and in the result add pairs with swapped elements to these pairs. Let's do it:

    If we have $x \ge y$, then $2x \ge 2y$, so $y^2+1 \ge x^2+1$, which gives $y^2 \ge x^2$. From this, it does \textit{not} follow in general that $y \ge x$, this holds only for positive $x$ and $y$. However, our $x$ and $y$ are certainly positive, for example in the first equation we have $x^2+1=2y$, on the left is a positive number, so $2y$ is also positive, so $y$ as well, and similarly $x$ too. Thus, from $y^2 \ge x^2$ it really follows that $y \ge x$. We thus have $x \ge y$ and $y \ge x$, which together gives $x=y$. We then finish the solution as before and obtain the single solution $(1,1)$.
}

Let us note that this system has no solution where the variables would not be equal, so we did not have to add others to the found solution $(1,1)$. In the following very simple example, it will not be so:

\Example{}{Solve the system in real numbers:
    $$
    \align
    |x-y| &= 2, \\
    x+y &=4. \\
    \endalign
    $$
}{
    The system is symmetric, so without loss of generality $x \ge y$. Thereby we get a system of two linear equations with two unknowns:
    $$
    \align
    x-y &= 2, \\
    x+y &=4, \\
    \endalign
    $$
    of which we can easily convince ourselves that it has the unique solution $(3,1)$. The original system therefore has two solutions, $(3,1)$ and $(1,3)$.
}

Symmetry is a fascinating thing. Its use by ordering variables is a very strong technique used in all areas of mathematics. When can we actually do this? In general: When by swapping variables $X$ and $Y$ we get the same problem, it suffices to solve the problem in the case $X \ge Y$ and argue that in the case $Y \ge X$ the problem is symmetric.

In the context of systems of equations, we might have a case where we have an equation with multiple variables, which is symmetric if and only if by swapping any two variables we get the same problem: In such a case we can order all variables, for example for a problem with variables $x,y,z$ it suffices to solve the case $x \ge y \ge z$ and permute the final solutions, for example if we had the solution $(3,2,1)$, then we have another 5 solutions $(3,1,2)$, $(2,1,3)$, $(2,3,1)$, $(1,3,2)$.

In the next example, we will show that we can order something even when we do not have complete symmetry:

\Example{}{
    Solve the system in real numbers:
    $$
    \align
    x^3 + 1 &=2y, \\
    y^3 + 1 &=2z, \\
    z^3 + 1 &=2x. \\ 
    \endalign
    $$
}{
    This system \textit{is not} symmetric: by swapping variables $x$ and $y$ we do not get exactly the same set of three equations. However, we see that it still exhibits a form of symmetry --- the equations somehow repeat the variables $x,y,z$ \uv{in a circle} in this order. More precisely: If we replace variables $(x,y,z)$ with variables $(z,x,y)$, we get the same system (verify this). We call such systems \textit{cyclic}. In them, we cannot afford to order all variables as in symmetric ones, however, we can afford without loss of generality to assume that one of the variables is the largest/smallest. If e.g. $x$ is the maximum of $x,y,z$ and we find a solution for example $(3,2,1)$, then thanks to the cyclicity of the system we obtain another two solutions $(2,1,3)$ and $(1,2,3)$.

    After a long introduction, let's finally get into the example. Since our system is cyclic, it suffices to solve the case where $x=\max\{x,y,z\}$. Then it holds $x \ge z$, thus $x^3+1 \ge z^3+1$ (since the function $f\colon x\rightarrow x^3$ is increasing). But from the comparison of the first and third equation this means that $2y \ge 2x$, or $y \ge x$. This together with $x \ge y$ gives $x=y$. This allows us to use $y \ge z$ in a similar way, $2z=y^3+1 \ge z^3+1 = 2x$, so $z \ge x$, thus $z=x$, altogether $x=y=z$.
    
    All equations therefore have the form $t^3+1=2t$, equivalently $t^3-2t+1=0$. Although this is a cubic equation, fortunately we can easily guess its root $t=1$, so by the standard method of polynomial division we arrive at the factorization:
    $$
    t^3-2t+1 =(t-1)(t^2+t-1)
    $$
    We solve the last equation playfully and find that it has three solutions $t=1$ and $t=\frac12(-1\pm\sqrt5)$, which thanks to the equivalence of the used adjustments also satisfy the original equation $t^3+1=2t$. 
    
    The system has three solutions of the form $(x,y,z)=(t,t,t)$, where $t \in \{1, \frac12(-1\pm\sqrt5)\}$.
}

\Exercise{}{
    In the domain of real numbers solve the system of equations
    $$
    \align
    x &= \sqrt{y+2},\\
    y &= \sqrt{x+2}.\\
    \endalign
    $$
}{
    The system is symmetric, without loss of generality let us assume that $x \ge y$. Since the function $f(t)=\sqrt{t+2}$ is increasing, it holds
    $$
    x \ge y \implies \sqrt{x+2} \ge \sqrt{y+2},
    $$
    so using the system we have $y \ge x$. From $x \ge y$ and $y \ge x$ it necessarily follows $x=y$.
    
    By substituting into the first equation we get $x=\sqrt{x+2}$. After squaring we have $x^2=x+2$, which is a quadratic equation with roots $x=2$ and $x=-1$, so $(x,y)$ is either $(2,2)$ or $(-1,-1)$. However, only the first pair passes the check.
    
    The system has a unique solution $(x,y)=(2,2)$.
}

\Exercise{}{
    Determine all triples $(x,y,z)$ of real numbers for which it holds
    $$
    \align
    (x+y)^5 &= 32z, \\
    (y+z)^5 &= 32x, \\
    (z+x)^5 &= 32y. \\
    \endalign
    $$
}{
    The system is symmetric, therefore without loss of generality let us assume that $x \ge y \ge z$. Then also $32x \ge 32y \ge 32z$, according to the system we therefore have
    $$
    (y+z)^5 \ge (z+x)^5 \ge (x+y)^5.
    $$
    Since the function $f\colon t \rightarrow t^5$ is increasing, we get from this
    $$
    y+z \ge z+x \ge x+y,
    $$
    resp. $y \ge x$ from the first and $z \ge y$ from the second. Together with the assumption $x \ge y \ge z$ this means $x=y=z$. Substituting into the first equation we get $(2x)^5=32x$.
    $$
    32x^5-32x=0 \implies x(x^5-1)=0,
    $$   
    so $x=0$ or $x=1$.

    We can convince ourselves by checking that the found solutions $(0,0,0)$ and $(1,1,1)$ indeed satisfy the system.
}

\Exercise{}{
    Determine all triples $(x,y,z)$ of real numbers for which it holds
    $$
    \align
    x^5 &= 5y^3 - 4z, \\
    y^5 &= 5z^3 - 4x, \\
    z^5 &= 5x^3 - 4y. \\
    \endalign
    $$
}{
    The system is not symmetric, so we cannot assume $x \ge y \ge z$. However, it is cyclic, without loss of generality let $x=\max\{x,y,z\}$. Thus $x \ge z$, and therefore $x^5 \ge z^5$. Let's look at what the first and third equations tell us:
    $$
    \align
    5y^3 - 4z &\ge 5x^3 - 4y, \\
    4y - 4z &\ge 5x^3 - 5y^3.
    \endalign
    $$
    Since however also $x \ge y$, then also $5x^3 \ge 5y^3$, so the last inequality means even $4y-4z \ge 0$, thus $y \ge z$. Thus $y^5 \ge z^5$. Let's look at what the second and third equations tell us:
    $$
    \align
    5z^3-4x &\ge 5x^3 - 4y\\
    4y - 4x &\ge 5x^3-5z^3.
    \endalign
    $$
    Similar to the previous comparison we have $y \ge x$. Together with $x \ge y$ this however means $x=y$. The last inequality then means $0 \ge 5x^3-5z^3$, so $z^3 \ge x^3$, which gives $z \ge x$, which together with $x \ge z$ means also $x=z$. Necessarily therefore $x=y=z$.

    Substituting $x=y=z$ into the first equation we get $x^5 = 5x^3 - 4x$, i.e., $x^5-5x^3+4x=0$. We see that $x=0$ is a root, thus we have $x(x^4-5x^2+4)=0$. The polynomial $x^4-5x^2+4$ is obviously quadratic in the variable $t=x^2$, by solving the equation $t^2-5t+4=0$ we easily find roots $t=4$ and $t=1$. Our equation can therefore be factored as $x(x^2-4)(x^2-1)=0$, so in the end it has 5 roots: $0$, $\pm 1$, $\pm 2$.

    We can convince ourselves by checking that all 5 found solutions $(0,0,0)$, $(-1,-1,-1)$, $(1,1,1)$, $(-2,-2,-2)$, $(2,2,2)$ satisfy the system.
}

\sec What to remember

\secc Solving Techniques

\begitems \style N
\i Add and subtract equations. In some problems, it may also help to multiply or divide them.
\i Look for factorizations of expressions. 
\i Observe symmetry or cyclicity and without loss of generality order the variables or assume which is the largest or smallest. Do not forget to include symmetric or cyclic solutions in the final answer.
\i Notice if expressions cannot be adjusted to sums of squares. 
\i Do not forget the check -- even if technically not necessary, often it is easier to do it than to precisely justify that it is not needed.
\enditems

\textit{Note.} There are many techniques which we did not get to -- the most common of them is the use of more complex inequalities. We will return to such systems in the material on inequalities.

\secc General useful knowledge

\begitems \style N
\i Whenever we see $x^2+y^2+z^2 = xy+yz+zx$, it immediately means $x=y=z$ (actually the left side is always at least as large as the right and equality occurs only for $x=y=z$)
\i Symmetry and cyclicity are not tied only to systems of equations -- in all parts of mathematics it is very important to assume something without loss of generality and thereby simplify the problem.
\enditems

\sec Further Problems

So that it's not too little, here are a few more problems. All combine in various ways the methods which we have learned or seen along the way (such as factorization, completing the square, case analysis, adding/subtracting equations, etc.). Problems are ordered approximately by difficulty.

\Problem{0}{74th year MO, regional round B, problem 1}{
    In the domain of real numbers solve the system of equations
    $$
    \align
    x^2+4y^2+z^2-4xy-2z+1&=0,\\
    y^2-xy-2y+2x&=0.
    \endalign
    $$
}{
    Notice that the second equation can actually be factored. This helps in simplifying the first equation.
}{
    The second equation gives us $(y-x)(y-2)=0$. This gives two cases, $y=x$ and $y=2$. Substitute into the first equation and let's notice that it looks very \textit{square-like}.
}{
    We factor the second equation by successive factoring out:
    $$
    y(y-2) - x(y-2) = 0 \implies (y-x)(y-2) = 0.
    $$
    This equality occurs if and only if $y=x$ or $y=2$. We will analyze both cases.

    \begitems
    \i If $y=x$, substituting into the first equation we get
    $$
    x^2+4x^2+z^2-4x^2-2z+1=0,
    $$
    which after adjustment changes to
    $$
    x^2 + (z-1)^2 = 0.
    $$
    Since this is a sum of two non-negative numbers, equality can occur only if both addends are zero. Thus $x=0$ and $z=1$. Since $y=x$, we have the solution $(0,0,1)$.
    \i If $y=2$, substituting into the first equation we get
    $$
    x^2+4(2^2)+z^2-4x(2)-2z+1 = 0,
    $$
    which after adjustment and completing the square gives
    $$
    \align
    x^2-8x+16+z^2-2z+1&=0,\\
    (x-4)^2 + (z-1)^2 &= 0.
    \endalign
    $$
    Similar to the first case, from here it follows $x=4$ and $z=1$. Since we assumed $y=2$, we have the solution $(4,2,1)$.
    \enditems
    
    By checking we convince ourselves that both found triples $(0,0,1)$ and $(4,2,1)$ are solutions of the system.
}


\Problem{0}{57th year MO, national round A in Czechia\fnote{In this year, the national and regional rounds were \textit{exceptionally} different in the Czech Republic and Slovakia.}, problem 1}{
    Solve the system in real numbers:
    $$
    \align
    x + y^2 &= y^3, \\
    y + x^2 &= x^3. \\
    \endalign
    $$
}{
    Subtract the equations and factor. 
}{
    After subtracting the equations and factoring we have that either $x=y$ or $x^2+xy+y^2-x-y+1=0$. The second equation is non-evident. A straightforward way to solve it is to consider one variable (e.g., $x$) as an unknown and the second (e.g., $y$) as a parameter and solve the quadratic equation. However, a nicer way is to find a factorization into a sum of squares, for which similar to the analysis of $x^2+y^2+z^2-xy-yz-zx$ we first need to multiply by two.
}{
    Subtracting the equations we get $x-y-(x^2-y^2)=y^3-x^3$, which after factoring out the term $(x-y)$ leads to
    $$
    (x-y)(1 - x - y + x^2+xy+y^2) = 0.
    $$
    This gives us two possibilities.

    \begitems
    \i If $x=y$, substituting into the first equation of the system we have $x+x^2=x^3$, i.e., $x(x^2-x-1)=0$.
    The solutions are $x=0$ and the roots of the equation $x^2-x-1=0$, thus $x=\frac{1\pm\sqrt{5}}{2}$.
    We thus obtain three solutions: $(0,0)$, $(\frac{1+\sqrt{5}}{2}, \frac{1+\sqrt{5}}{2})$ and $(\frac{1-\sqrt{5}}{2}, \frac{1-\sqrt{5}}{2})$.

    \i If $x^2+xy+y^2-x-y+1=0$, we multiply the equation by two and appropriately rearrange the terms into a sum of three squares:
    $$
    \align
    2x^2+2xy+2y^2-2x-2y+2&=0, \\
    (x^2-2x+1) + (y^2-2y+1) + (x^2+2xy+y^2) &= 0, \\
    (x-1)^2 + (y-1)^2 + (x+y)^2 &= 0.
    \endalign
    $$
    The sum of three non-negative numbers is zero if and only if all three are zero. It would therefore have to hold $x-1=0$, $y-1=0$ and $x+y=0$. The first two equations give $x=1$ and $y=1$, which is in contradiction with the third equation. This case therefore has no solution.
    \enditems

    The system therefore has exactly the three solutions listed in the first case.
}

\Problem{0}{70th year MO, school round B, problem 1}{
    For real numbers $x$, $y$, $z$ it holds
    $$
    \align
    |x+y|&=1-z,\\
    |y+z|&=1-x,\\
    |z+x|&=1-y.
    \endalign
    $$
    Find out what values the sum $x+y+z$ can take. For each complying sum, give an example of corresponding numbers $x$, $y$, $z$.
}{
    Squaring works well with absolute values, thanks to that they disappear.
}{
    After removing absolute values, it is possible for example to subtract each pair of equations and factor (our goal is not to find all solutions, but only possible values of the sum $x+y+z$, which will help us in the discussion). An even faster solution is surprisingly to add all equations. In any case, do not forget to verify that for the found candidates for $x+y+z$ there indeed exist $x,y,z$ for which the sum is attainable.
}{
    To get rid of absolute values, we square each of the three equations:
    $$
    \align
    (x+y)^2 &= (1-z)^2, \\
    (y+z)^2 &= (1-x)^2, \\
    (z+x)^2 &= (1-y)^2.
    \endalign
    $$
    Adding these three equations and expanding all brackets we get:
    $$
    2x^2 + 2y^2 + 2z^2 + 2xy + 2yz + 2zx = 3 - 2(x+y+z) + x^2 + y^2 + z^2.
    $$
    We simplify this equation by subtracting the expression $x^2+y^2+z^2$ from both sides. We are left with:
    $$
    x^2+y^2+z^2 + 2xy+2yz+2zx = 3 - 2(x+y+z).
    $$
    The expression on the left side is exactly the formula for $(x+y+z)^2$. If we denote the sum $S = x+y+z$, the equation transforms to $S^2 = 3 - 2S$, i.e., $S^2 + 2S - 3 = 0$. By factoring $(S+3)(S-1)=0$ we find two possible values for the sum: $S=1$ and $S=-3$.

    We still must verify if both these values are truly attainable.
    \begitems
    \i For $S=1$ we can choose $(x,y,z)=(1,0,0)$. We easily verify that this satisfies the system.
    \i For $S=-3$ it suffices to take $(x,y,z)=(-1,-1,-1)$.
    \enditems
    Possible values of the sum $x+y+z$ are therefore $1$ and $-3$.
    
    Let us add that the problem has other possible ways of solution, we can for example discuss signs of expressions $x+y$, $y+z$, $z+x$. 
}

\Problem{0}{71st year MO, regional round A, problem 2}{
    In the domain of \textit{positive\fnote{A not insignificant number of participants did not notice the positivity and solved a problem that did not have nice solutions.}} real numbers solve the system of equations
    $$
    \align
    x^2 + 2y^2 &= \hphantom{1}x + 2y + 3z, \\
    y^2 + 2z^2 &= 2x + 3y + 4z, \\
    z^2 + 2x^2 &= 3x + 4y + 5z.
    \endalign
    $$
}{
    Notice that after adding all equations we have on the right side $6x+9y+12z$, which is somewhat suspiciously similar to the right side of one of our equations.
}{
    The trick is to add all equations and subtract twice the second equation, thereby completely zeroing out the right side and leaving $x^2=z^2$ on the left, which thanks to the positivity of $x,y,z$ gives $x=z$.
}{
    Adding all three equations we get
    $$
    3x^2+3y^2+3z^2 = 6x+9y+12z.
    $$
    The right side is exactly three times the right side of the second equation. Subtracting three times the second equation from the sum of all equations we thus obtain an equation with a zero right side:
    $$
    (3x^2+3y^2+3z^2) - 3(y^2+2z^2) = 0.
    $$
    After adjustment we have $3x^2 - 3z^2 = 0$, i.e., $x^2=z^2$. Since $x,y,z$ are positive real numbers, it follows from this that $x=z$.

    Substituting $z=x$ into the first two equations of the system we get:
    $$
    \align
    x^2 + 2y^2 &= 4x + 2y = 2(2x+y), \\
    y^2 + 2x^2 &= 6x + 3y = 3(2x+y), \\
    \endalign
    $$
    We see that $3(x^2+2y^2) = 2(y^2+2x^2)$, from which $4y^2=x^2$, thus from positivity we have $2y=x$. By backward substitution into e.g. the first equation we get $(2y)^2 + 2y^2 = 2(4y+y)$, thus $6y^2=10y$, so either $y=0$ or $y=5/3$. 
    
    Altogether we have two candidates for solutions: $(0,0,0)$, $(10/3,5/3,10/3)$. By checking we easily verify that both triples satisfy the system.
}

\Problem{0}{}{
    In the domain of real numbers solve the system of equations
    $$
    \align
    a^3+b &= c, \\
    b^3+c &= d, \\
    c^3+d &= a, \\
    d^3+a &= b.
    \endalign
    $$
}{
    The system is cyclic, so it suffices to solve the case when $a=\max\{a,b,c,d\}$. Thus $a^3 \ge b^3$ etc., we have more options how to use it. One of them gives us another inequality between $a,b,c,d$.
}{
    The key is to use $a^3 \ge c^3$, this gives us $c-b \ge a-d$, which after adjustment means $d-b \ge a-c$, so necessarily $d \ge b$. What does $d^3 \ge b^3$ give us?
}{
    The system is cyclic, therefore it suffices to solve $a=\max\{a,b,c,d\}$. Then $a \ge c$, so also $a^3 \ge c^3$. From the first and third equation we thus get $c-b \ge a-d$, which we adjust to $d-b \ge a-c$. Since $a \ge c$, the right side is non-negative, thus $d \ge b$.

    From the inequality $d \ge b$ we similarly derive $d^3 \ge b^3$, which after substitution from the second and fourth equation gives $b-a \ge d-c$, thus $b-d \ge a-c$. But $a \ge c$, thus $b \ge d$, which together with $d \ge b$ gives $d=b$, and subsequently $b-d \ge a-c$ gives $0 \ge a-c$, so $c=a$.
    
    By substituting into the system we reduce it to
    $$
    \align
    a^3+b &= a, \\
    b^3+a &= b.
    \endalign
    $$
    Adding the equations we get $a^3+b^3=0$, so $a^3=(-b)^3$, from which $b=-a$. By substituting we now have only one equation $a^3-a=a$, which is equivalent to $a(a^2-2)=0$, so $a=0$ or $a=\pm \sqrt2$.

    By checking we easily verify that all quadruples
    $$
    (0,0,0,0), 
    (\sqrt{2}, -\sqrt{2}, \sqrt{2}, -\sqrt{2}), 
    (-\sqrt{2}, \sqrt{2}, -\sqrt{2}, \sqrt{2})
    $$
    are a solution.
}

\Problem{0}{}{
    Find the smallest four-digit natural number $n$, for which the system
    $$
    \align
    x^3+y^3+x^2y+y^2x&=n,\\
    x^2+y^2+x+y&=n+1
    \endalign
    $$
    has only integer solutions.
}{
    The first step is to notice that the left side of the equation can actually be factored into a product of expressions suspiciously similar to those in the second equation.
}{
    It holds $x^3+y^3+x^2y+y^2x=(x+y)(x^2+y^2)$. The second equation is meanwhile $(x+y)+(x^2+y^2)=n+1$. The numbers $A=x+y$, $B=x^2+y^2$ therefore satisfy $AB=n$, $A+B=n+1$. Which ones can they be? Formally we can derive them either by realizing that from Vieta's\fnote{Vieta's formulas are named after the French mathematician Francois Viète (1540--1603), who is considered one of the fathers of modern algebra thanks to his systematic use of letters to represent unknowns and parameters, which allowed solving problems more generally.} formulas $A,B$ must be roots of the quadratic equation $t^2-(n+1)t+n=0$ and solving it; or we can substitute $B=n+1-A$ into the first equation, etc. In any case we obtain two possibilities for $(A,B)$. One of them can be excluded quickly (let's not forget that $n$ is a four-digit number, thus it is positive).
}{
    It holds $x^3+y^3+x^2y+y^2x=(x+y)(x^2+y^2)$. If we denote $A=x+y$ and $B=x^2+y^2$, the system can be equivalently rewritten to the form $AB=n$ and $A+B=n+1$. From Vieta's formulas it follows that $A$ and $B$ are roots of the quadratic equation $t^2-(n+1)t+n=0$, which we can factor as $(t-1)(t-n)=0$. Possible values for the pair $\{A,B\}$ are therefore $\{1,n\}$. We will analyze both cases.

    \begitems
    \i Let $x+y=n$ and $x^2+y^2=1$. Since $x,y$ are integers, from the second equation we have $x,y \leq 1$. The sum $x+y$ equal to $n$ will therefore certainly not be a four-digit number.
    \i Let $x+y=1$ and $x^2+y^2=n$. At least one of the numbers $x,y$ is positive (otherwise $x+y \leq 0$), without loss of generality let it be $x$. 
    
    Substituting $y=1-x$ into the second equation we get $x^2+(1-x)^2=n$, thus $n=2x^2-2x+1$. The function $f(x)=2x^2-2x+1$ is obviously increasing for $x \ge 1$. The goal is to find the smallest $x$ for which $f(x)$ is four-digit. This can be done by solving the inequality $2x^2-2x+1 \ge 1000$, or by trying, the answer is $x=23$, since $f(22)=925$ and $f(23)=1013$. 
    
    On the other hand, for $n=1013$ the equation $2x^2-2x+1=1013$ has two solutions $x=23$ and $x=-22$, to which correspond $y=-22$ and $y=23$, so both are integer.

    The answer is $n=1013$.
    \enditems
}

\Problem{0}{64th year MO, national round A, problem 4}{
    In the domain of real numbers solve the system of equations
    $$
    \align
    a(b^2 + c) &= c(c + ab),\\
    b(c^2 + a) &= a(a + bc),\\
    c(a^2 + b) &= b(b + ca).
    \endalign
    $$
}{
    Adding and subtracting leads to nothing factorable. The trick in this problem is to use \textit{multiplication} of equations. For example, if we multiply the equations as they are, we can cancel in the case $abc \neq 0$. This gives us something, but I will reveal that nothing extra useful. The key is to multiply the equations after a minor adjustment of each of them.
}{
    We adjust the first equation by rearranging terms:
    $$
    \align
    a(b^2+c) &= c(c+ab), \\
    ab^2 + ac &= c^2 + abc, \\
    ab^2 - abc &= c^2 - ac, \\
    ab(b-c) &= c(c-a).
    \endalign
    $$
    Now we can multiply and assuming mutual distinctness of $a,b,c$ and their non-zero status it yields useful things. It will not be the end, but we will be close.
}{
    We adjust the first equation by rearranging terms:
    $$
    \align
    a(b^2+c) &= c(c+ab), \\
    ab^2 + ac &= c^2 + abc, \\
    ab^2 - abc &= c^2 - ac, \\
    ab(b-c) &= c(c-a).
    \endalign
    $$
    When we adjust the remaining equations in this way, we get:
    $$
    \align
    ab(b-c) &= c(c-a), \tag1 \\
    bc(c-a) &= a(a-b), \tag2 \\
    ca(a-b) &= b(b-c). \tag3
    \endalign
    $$
    We will analyze several cases.

    \begitems
    \i If any variable is zero, for example $a=0$, then from the first equation we have $0=c^2$, so $c=0$, which then in the third equation gives $b=0$. We have $(0,0,0)$, which is indeed a solution. Let us assume therefore that $a,b,c$ are non-zero.
    \i If all variables are non-zero, but some two are equal, for example $a=b$, then from the second equation $bc(c-a)=0$ gives thanks to non-zero $b,c$ the equality $c=a$, so all three variables are equal. We easily verify that every triple $(t,t,t)$ is a solution of the system (it includes already the found solution $(0,0,0)$). Let us assume therefore that $a,b,c$ are pairwise distinct. Multiplying equations (1), (2), (3) we get
    $$
    a^2b^2c^2(a-b)(b-c)(c-a) = abc(a-b)(b-c)(c-a).
    $$
    Since the variables are distinct and non-zero, we can divide by the expression $abc(a-b)(b-c)(c-a)$ and we get $abc=1$. We use this to adjust equations (1), (2), (3) by multiplying them successively by $c$, $a$, $b$:
    $$
    \align
    (b-c) &= c^2(c-a), \\
    (c-a) &= a^2(a-b), \\
    (a-b) &= b^2(b-c). 
    \endalign
    $$
    This will lead to a contradiction: The system is cyclic, so it suffices to solve $a=\max\{a,b,c\}$. Then $c-a<0$, so from the first equation $b-c<0$, so from the last equation $a-b<0$, contradiction.
    \enditems

    Summing up all cases we obtain that the only solutions are triples $(t,t,t)$ for any real number $t$.
}

\DisplayExerciseSolutions

\DisplayHints

\DisplayProblemSolutions

\bye